% Preámbulo del Libro de Ingeniería (LuaLaTeX). Generado por build_engineering_notebook.py desde templates/.
\documentclass[11pt,a4paper,oneside]{report}
\usepackage[a4paper,margin=2.2cm]{geometry}
\usepackage{fontspec}
\directlua{luaotfload.add_fallback("nbfallback", {"DejaVu Sans:mode=harf", "DejaVu Sans Mono:mode=harf"})}
\setmainfont{Latin Modern Roman}[RawFeature={fallback=nbfallback}]
\setsansfont{Latin Modern Sans}[RawFeature={fallback=nbfallback}]
\setmonofont{DejaVu Sans Mono}[Scale=0.82]
\usepackage[spanish,es-noshorthands]{babel}
\usepackage{graphicx}
\setkeys{Gin}{width=\linewidth,height=0.42\textheight,keepaspectratio}
\usepackage{float}
\usepackage{xcolor}
\usepackage{array}
\usepackage{calc}
\usepackage{longtable}
\usepackage{booktabs}
\usepackage{amssymb}
\providecommand{\sout}[1]{#1}
\usepackage{etoolbox}
\usepackage{fvextra}
\usepackage{listings}
\usepackage{tikz}
\usetikzlibrary{arrows.meta,positioning,shapes.geometric}
\usepackage{fancyhdr}
\usepackage{titlesec}
\usepackage[hidelinks,unicode]{hyperref}
\usepackage{xurl}
\usepackage{microtype}

\definecolor{nbblue}{HTML}{1F4E79}
\definecolor{nbgray}{HTML}{666666}
\titleformat{\chapter}[display]{\normalfont\bfseries\color{nbblue}}{\Large\chaptertitlename\ \thechapter}{4pt}{\Huge}
\titlespacing*{\chapter}{0pt}{-10pt}{24pt}
\pagestyle{fancy}
\fancyhf{}
\fancyhead[L]{\small\color{nbgray}Robotica76 · Libro de Ingeniería · MakeX 2026 StormBreaker}
\fancyfoot[C]{\thepage}
\renewcommand{\headrulewidth}{0.3pt}
\fancypagestyle{plain}{\fancyhf{}\fancyfoot[C]{\thepage}\renewcommand{\headrulewidth}{0pt}}

\setlength{\parindent}{0pt}
\setlength{\parskip}{5pt plus 1pt}
\setlength{\emergencystretch}{3em}
\setlength{\LTleft}{0pt}
\setlength{\LTright}{0pt}
\setcounter{secnumdepth}{2}
\setcounter{tocdepth}{1}

% Tablas de pandoc más compactas
\AtBeginEnvironment{longtable}{\footnotesize\setlength{\tabcolsep}{3pt}\renewcommand{\arraystretch}{1.15}}
\providecommand{\tightlist}{\setlength{\itemsep}{0pt}\setlength{\parskip}{0pt}}

% Bloques de código: se parten líneas largas
\fvset{breaklines=true,breakanywhere=true,fontsize=\scriptsize,frame=lines,framesep=2mm}

% Listados (capítulo Software generado desde Doxygen)
\lstset{language=Python,basicstyle=\ttfamily\scriptsize,numbers=left,numberstyle=\tiny,
  breaklines=true,frame=single,keywordstyle=\color{blue},commentstyle=\color{gray},
  literate={á}{{\'a}}1 {é}{{\'e}}1 {í}{{\'i}}1 {ó}{{\'o}}1 {ú}{{\'u}}1 {ñ}{{\~n}}1
  {Á}{{\'A}}1 {É}{{\'E}}1 {Í}{{\'I}}1 {Ó}{{\'O}}1 {Ú}{{\'U}}1 {Ñ}{{\~N}}1
  {¿}{{?`}}1 {¡}{{!`}}1}

\newcommand{\nbsource}[1]{{\par\footnotesize\color{nbgray}Fuente: \texttt{#1}\par}}
